\section{The Kuga--Satake route and where it stops}\label{sec:ks}

\Cref{thm:hyperplane} closes induction on dimension through hyperplane
sections. One classical construction does reach a middle degree Hodge class
without any such induction, and we set it out in full, both because
it shows what a successful route looks like and because its boundary is
numerical and can be located exactly. The construction is Kuga and Satake's
\cite{KS67}, in the form given by Deligne \cite{Del72}. This section builds it,
proves that \Cref{thm:hyperplane} has no bearing on it, and then shows that for
abelian varieties of Weil type it is available precisely when $n=1$, which is
the one case in which the Weil line consists of divisor classes and there is
nothing to prove.

The arithmetic here is checked in \texttt{kugasatake.py}, item (VIII) of
\Cref{sec:verification}.%
\footnote{Kuga and Satake's paper is four pages and constructs the abelian variety; the Hodge-theoretic properties that make it useful, in particular the embedding of \Cref{prop:ksembed}, are Deligne's, and Rapoport's complement to that paper supplies the comparison with the Shimura variety. Satake had already set up the Clifford algebra families in \cite{Sat66}, which is where the Riemann form comes from.}

\subsection{The construction}

\begin{setupenv}[K3 type]\label{setup:k3type}
Let $T$ be a $\QQ$-vector space of dimension $b$ with a nondegenerate
symmetric form $q$, and let $T$ carry a Hodge structure of weight two,
$T_{\CC}=T^{2,0}\oplus T^{1,1}\oplus T^{0,2}$, polarised by $q$, with
\[
  \dim_{\CC}T^{2,0}=1 .
\]
Such a structure is said to be \emph{of K3 type}. Write $T^{2,0}=\CC\sigma$;
the Hodge--Riemann relations give $q(\sigma,\sigma)=0$ and
$q(\sigma,\bbar{\sigma})>0$, so the real plane
$P=\langle\Re\sigma,\Im\sigma\rangle\subset T_{\RR}$ is positive definite and
$q$ has signature $(2,b-2)$ on $T_{\RR}$. Normalise $\sigma$ so that
$e_{1}=\Re\sigma$ and $e_{2}=\Im\sigma$ are orthonormal.
\end{setupenv}

\begin{definition}[Clifford algebra]\label{def:clifford}
$C(T,q)$ is the quotient of the tensor algebra on $T$ by the ideal generated
by $v\otimes v-q(v)$ for $v\in T$. It is $\ZZ/2$-graded, $C=C^{+}\oplus C^{-}$,
by the parity of the word length.
\end{definition}

\begin{lemma}\label{lem:cliffdim}
For an orthogonal basis $f_{1},\dots,f_{b}$ of $T$ the products $f_{S}$ over
subsets $S\subseteq\{1,\dots,b\}$, taken in increasing order, form a basis of
$C(T,q)$. Hence $\dim_{\QQ}C=2^{b}$ and $\dim_{\QQ}C^{\pm}=2^{b-1}$, and
$C^{-}\cdot C^{+}\cdot C^{-}\subseteq C^{+}$.
\end{lemma}

\begin{proof}
The relations $f_{i}^{2}=q(f_{i})$ and $f_{i}f_{j}=-f_{j}f_{i}$ for $i\ne j$,
the latter obtained by expanding $(f_{i}+f_{j})^{2}=q(f_{i})+q(f_{j})$, let one
rewrite any word in the $f_{i}$ as a scalar times an increasing product, so the
$f_{S}$ span; they are independent by the standard filtration argument
comparing $C$ with $\bigwedge^{*}T$. The dimension counts are
$\sum_{k}\binom{b}{k}=2^{b}$ and
$\sum_{k\ \mathrm{even}}\binom{b}{k}=2^{b-1}$. The last statement is the
grading: odd plus even plus odd is even.
\end{proof}

\begin{lemma}[The complex structure]\label{lem:cliffJ}
With $e_{1},e_{2}$ as in \Cref{setup:k3type}, put $J=e_{1}e_{2}\in C^{+}$.
Then $J^{2}=-1$, and left multiplication by $J$ is an $\RR$-linear operator on
$C^{+}\otimes_{\QQ}\RR$ whose square is $-\id$.
\end{lemma}

\begin{proof}
$J^{2}=e_{1}e_{2}e_{1}e_{2}=-e_{1}e_{1}e_{2}e_{2}=-q(e_{1})q(e_{2})=-1$, using
anticommutativity of $e_{1}$ and $e_{2}$ and then $e_{i}^{2}=q(e_{i})=1$.
Left multiplication by an element of $C^{+}$ preserves $C^{+}$, and the square
of left multiplication by $J$ is left multiplication by $J^{2}=-1$.
\end{proof}

\begin{definition}[Kuga--Satake variety]\label{def:ksvariety}
Fix a lattice $C^{+}_{\ZZ}\subset C^{+}$ stable under multiplication. The
complex torus
\[
  \mathrm{KS}(T)=\bigl(C^{+}\otimes_{\QQ}\RR,\,J\bigr)\big/\,C^{+}_{\ZZ}
\]
has dimension $\tfrac12\dim_{\QQ}C^{+}=2^{b-2}$. Satake's construction of a
Riemann form on $C^{+}$ compatible with $J$ \cite{Sat66} makes it an abelian
variety, and by construction
\[
  H^{1}\bigl(\mathrm{KS}(T),\QQ\bigr)=C^{+}(T)
\]
as a Hodge structure of weight one, the Hodge decomposition being the
eigenspace decomposition of left multiplication by $J$.
\end{definition}

\begin{example}\label{ex:ksdims}
For a projective K3 surface, $T$ is the transcendental lattice inside
$H^{2}(X,\QQ)$; taking all of $H^{2}$ with $b=22$ gives
$\dim\mathrm{KS}=2^{20}=1{,}048{,}576$. At the other end, $b=3$ gives
$\dim\mathrm{KS}=2$, an abelian surface; this is the case relevant below.
\end{example}

\subsection{What the construction transports}

\begin{proposition}[The embedding]\label{prop:ksembed}
Fix $v_{0}\in T$ with $q(v_{0})\ne0$. The map
\[
  \Psi\colon T\longrightarrow\End_{\QQ}\bigl(C^{+}\bigr),
  \qquad
  \Psi(v)(x)=v\,x\,v_{0},
\]
is injective and is a morphism of Hodge structures
$T(1)\to\End\bigl(H^{1}(\mathrm{KS}(T),\QQ)\bigr)$, both sides having weight
zero. Consequently a Hodge class in $T$, that is a rational class of type
$(1,1)$, is carried to a Hodge class of type $(0,0)$ in
$\End(H^{1}(\mathrm{KS}(T),\QQ))$, that is to an element of
$\End^{0}(\mathrm{KS}(T))$.
\end{proposition}

\begin{proof}
The map lands in $\End(C^{+})$ by the last part of \Cref{lem:cliffdim}, since
$v,v_{0}\in T\subset C^{-}$. It is injective because
$\Psi(v)(v_{0}^{-1})=v$, where $v_{0}^{-1}=v_{0}/q(v_{0})$. That it respects
the Hodge structures is Deligne's computation \cite[\S4]{Del72}: the
Hodge structure on $\End(C^{+})$ is induced by $J$ acting on both sides, and
$\Psi$ intertwines the weight-two structure on $T$, twisted once, with it.
The last sentence is the definition of a morphism of Hodge structures together
with the theorem of Lefschetz on $(1,1)$ classes in the form that says a Hodge
class of type $(0,0)$ in $\End(H^{1}(B,\QQ))$ of an abelian variety $B$ is an
element of $\End^{0}(B)$.
\end{proof}

\begin{remark}[What is and is not known about the transport]%
\label{rem:kstransport}
The graph of an element of $\End^{0}(\mathrm{KS}(T))$ is an algebraic cycle on
$\mathrm{KS}(T)\times\mathrm{KS}(T)$, so the image of a Hodge class under
$\Psi$ is algebraic. Carrying that conclusion back to the variety $X$ from
which $T$ came requires the correspondence $\Psi$ itself to be induced by an
algebraic cycle on $X\times\mathrm{KS}(T)\times\mathrm{KS}(T)$, and that is
what is not known in general. What is known is that $\Psi$ is absolute Hodge,
by Deligne's theorem for abelian varieties together with the analysis of the
Kuga--Satake correspondence, and that it is motivated in Andr\'e's sense
\cite{And96}; van Geemen \cite{vG00} surveys the cases where algebraicity has
been established. So the Kuga--Satake route converts the Hodge conjecture for
$T$ into the algebraicity of one specific correspondence, which is progress of
a definite kind and is not a proof.
\end{remark}

\subsection{The hyperplane obstruction does not apply}

\begin{proposition}\label{prop:ksnotrestriction}
\Cref{thm:hyperplane,thm:returntrip} do not constrain the Kuga--Satake route.
\end{proposition}

\begin{proof}
Both theorems are statements about the pair of maps $j^{*}$ and $j_{*}$
attached to the inclusion $j\colon Y\hookrightarrow X$ of a smooth member of an
ample linear system, and both proofs use two properties of that pair which no
other correspondence has: the Lefschetz hyperplane theorem, which makes
$j^{*}$ an isomorphism below the middle degree and hence forces the
vanishing of \Cref{lem:gysinrestrict} in the primitive middle degree; and the
identity $j_{*}j^{*}=\,\cdot\,\eta$, which is what makes the return trip
multiply by the polarisation rather than recover the class. The Kuga--Satake
correspondence is a cycle class on a triple product of varieties of unrelated
dimensions, $\dim\mathrm{KS}(T)=2^{b-2}$ against $\dim X$, and it satisfies
neither property; in particular the composite of $\Psi$ with its transpose is
multiplication by $q(v_{0})$ on $T$ and not by a power of any polarisation of
$X$. So neither hypothesis of \Cref{thm:hyperplane} is met.
\end{proof}

This is the content already asserted in \Cref{rem:ksscope}. The point of
isolating it is that it identifies what an escape from \Cref{thm:hyperplane}
must look like: a correspondence whose transpose composite is not the
Lefschetz operator. \Cref{sec:escape} makes the analogous identification for
\Cref{thm:divisor,thm:character}.

\subsection{Where the route stops for Weil type}

The limitation of the Kuga--Satake construction is the hypothesis
$\dim T^{2,0}=1$, and for abelian varieties of Weil type the possible values of
$\dim T^{2,0}$ can be listed completely.

\begin{proposition}[The pieces of \texorpdfstring{$H^{2}$}{H2}]\label{prop:h2pieces}
Let $(A,\iota,\eta)$ be of $(K,-1,n)$-Weil type. In the character grading of
\Cref{prop:grading},
\begin{equation}\label{eq:h2split}
  H^{2}(A,\CC)
  =\textstyle\bigwedge^{2}V_{+}\;\oplus\;\bigl(V_{+}\otimes V_{-}\bigr)
   \;\oplus\;\bigwedge^{2}V_{-},
\end{equation}
the middle summand is the isotypic component of the norm character and is
defined over $\QQ$, and the outer two are interchanged by conjugation so their
sum is defined over $\QQ$. The Hodge numbers are
\[
  \bigl(h^{2,0},h^{1,1},h^{0,2}\bigr)=
  \begin{cases}
    \bigl(\tbinom{n}{2},\,n^{2},\,\tbinom{n}{2}\bigr)
      & \text{on }\bigwedge^{2}V_{\pm},\\[4pt]
    \bigl(n^{2},\,2n^{2},\,n^{2}\bigr)
      & \text{on }V_{+}\otimes V_{-},
  \end{cases}
\]
and the total is $2\tbinom{2n}{2}+4n^{2}=\tbinom{4n}{2}=\dim H^{2}(A,\QQ)$.
\end{proposition}

\begin{proof}
The splitting is \eqref{eq:grading} at $k=2$. The characters are
$\tau^{2}$, $\tau\bbar{\tau}=\Nm(\tau)$ and $\bbar{\tau}^{2}$, so the middle
summand is the norm-isotypic component, and it is self-conjugate, hence defined
over $\QQ$; the outer two are exchanged. For the Hodge numbers use
\eqref{eq:bisplit} and \Cref{lem:balanced}: $V_{\pm}^{1,0}$ and
$V_{\pm}^{0,1}$ have dimension $n$ each, so
$(\bigwedge^{2}V_{+})^{2,0}=\bigwedge^{2}V^{1,0}_{+}$ has dimension
$\binom{n}{2}$ and $(\bigwedge^{2}V_{+})^{1,1}=V^{1,0}_{+}\otimes V^{0,1}_{+}$
has dimension $n^{2}$, while
$(V_{+}\otimes V_{-})^{2,0}=V^{1,0}_{+}\otimes V^{1,0}_{-}$ has dimension
$n^{2}$ and the $(1,1)$ part has dimension $2n^{2}$. The total is
$2(2\binom{n}{2}+n^{2})+4n^{2}=2(2n^{2}-n)+4n^{2}=8n^{2}-2n=\binom{4n}{2}$.
\end{proof}

\begin{proposition}[No K3-type piece beyond \texorpdfstring{$n=1$}{n=1}]%
\label{prop:normpart}
Let $(A,\iota,\eta)$ be of $(K,-1,n)$-Weil type with $\Hg(A)=\SU(V,H)$. Then
$H^{2}(A,\QQ)$ decomposes into exactly three rational sub-Hodge structures that
cannot be decomposed further, namely
\[
  \QQ\,\eta,
  \qquad
  \bigl(V_{+}\otimes V_{-}\bigr)_{\QQ}\ominus\QQ\eta,
  \qquad
  \Bigl(\textstyle\bigwedge^{2}V_{+}\oplus\bigwedge^{2}V_{-}\Bigr)_{\QQ},
\]
with $h^{2,0}$ equal to $0$, $n^{2}$ and $n(n-1)$ respectively. Consequently
$H^{2}(A,\QQ)$ contains a rational sub-Hodge structure of K3 type if and only
if $n=1$.
\end{proposition}

\begin{proof}
A rational sub-Hodge structure $W\subseteq H^{2}(A,\QQ)$ is stable under
$\Hg(A)$: the orthogonal projector onto $W$ for the polarisation is a Hodge
class in $\End(H^{2})$, hence is fixed by $\Hg(A)$, hence $W$ is a
$\Hg(A)$-submodule. Conversely a $\Hg(A)$-submodule is a sub-Hodge structure,
because the Hodge decomposition is the eigenspace decomposition of an element
of $\Hg(A)(\RR)$. So the rational sub-Hodge structures are the
$\QQ$-subrepresentations of $\SU(V,H)$.

After complexification $\SU(V,H)$ becomes $\SL(V_{+})$, acting on $V_{+}$
tautologically and on $V_{-}\cong V_{+}^{\vee}$ dually. In \eqref{eq:h2split}
the outer summands $\bigwedge^{2}V_{+}$ and $\bigwedge^{2}V_{-}$ are
irreducible, being fundamental representations, and the middle summand is
$V_{+}\otimes V_{+}^{\vee}=\End(V_{+})=\mathfrak{sl}(V_{+})\oplus\CC$, the sum
of the adjoint representation and the trivial one. This gives four complex
irreducible pieces; conjugation fixes the two coming from the middle and
exchanges the outer two, so over $\QQ$ they group into the three listed
structures, and no further splitting is possible.

The trivial summand is spanned by $\eta$ by \Cref{lem:NSchar}, and has
$h^{2,0}=0$ since $\eta$ is of type $(1,1)$. The adjoint summand has
$h^{2,0}=n^{2}$, since removing the one-dimensional trivial piece from
$V_{+}\otimes V_{-}$ removes a $(1,1)$ class only. The outer structure has
$h^{2,0}=2\binom{n}{2}=n(n-1)$ by \Cref{prop:h2pieces}.

For a sum of these to have $h^{2,0}=1$ we need one of $n^{2}$, $n(n-1)$,
$n^{2}+n(n-1)=2n^{2}-n$ to equal $1$, using $h^{2,0}(\QQ\eta)=0$. The equation
$n(n-1)=1$ has no positive integer solution, and each of $n^{2}=1$ and
$2n^{2}-n=1$ gives $n=1$. At $n=1$ the adjoint summand indeed has
$h^{2,0}=1$ and the outer one is zero.
\end{proof}

\begin{corollary}[The route is available exactly where it is not needed]%
\label{cor:ksboundary}
Let $(A,\iota,\eta)$ be of $(K,-1,n)$-Weil type with $\Hg(A)=\SU(V,H)$. The
Kuga--Satake construction applies to a rational sub-Hodge structure of
$H^{2}(A,\QQ)$ if and only if $n=1$; and for $n=1$ the Weil line lies in
$H^{2}(A,\QQ)$ and consists of divisor classes, hence of algebraic classes by
the Lefschetz theorem on $(1,1)$ classes, so nothing needs transporting.
\end{corollary}

\begin{proof}
The first clause is \Cref{prop:normpart} with \Cref{setup:k3type}. For the
second, $2n=2$, so $\HW(A)\subseteq H^{2}(A,\QQ)\cap H^{1,1}(A)$, which is
$\NS(A)\otimes\QQ$.
\end{proof}

\begin{remark}[Why there is no higher analogue on the shelf]%
\label{rem:nohigherks}
For $n\ge2$ the Weil line sits in $H^{2n}(A,\QQ)$ with $2n\ge4$, and the
Clifford construction has no counterpart in weight greater than two: it turns
a quadratic form into an algebra, and a Hodge structure of weight $w$ carries a
polarisation that is symmetric only for $w$ even and has no reason to produce
a Clifford algebra with the right complex structure. One might instead apply
the construction to a sub-Hodge structure of $H^{2}$ and hope to recover the
degree $2n$ class from the resulting abelian variety;
\Cref{prop:normpart} rules that out by showing the hypothesis fails for every
$n\ge2$. What remains conceivable is a construction of Kuga--Satake shape for
the weight $2n$ structure directly, producing an auxiliary variety on which
the Weil line becomes a divisor class or an endomorphism. We do not know of
one, and we note that \Cref{thm:character} does not obstruct it: such a
construction is a correspondence and need not produce a $\Kx$-eigenclass. This
is the same gap as in \Cref{rem:nonequivariant}, seen from the other side.
\end{remark}

\begin{figure}[htbp]
\centering
  \includegraphics[width=0.88\textwidth]{fig_lightcone}
\caption{The geometry the Kuga--Satake construction starts from, in the
smallest case $b=3$. The quadratic form is $q=x^{2}+y^{2}-z^{2}$ of signature
$(2,1)$, with its null cone; the horizontal plane $P=\{z=0\}$ is positive
definite and meets the cone only at the origin. The construction needs such a
plane, and it gets one from \Cref{setup:k3type}: $P=\langle\Re\sigma,
\Im\sigma\rangle$ for the period $\sigma$ spanning $T^{2,0}$. An orthonormal
basis $e_{1},e_{2}$ of $P$ gives $J=e_{1}e_{2}$ with $J^{2}=-1$, by
\Cref{lem:cliffJ}, and $J$ is the complex structure on
$\mathrm{KS}(T)$. Positive definite planes correspond to the points of one
sheet of the hyperboloid $q=-1$, which is the period domain. The hypothesis
that makes all of this available is $\dim T^{2,0}=1$, and by
\Cref{prop:normpart} an abelian variety of Weil type supplies it only when
$n=1$.}
\label{fig:ks}
\end{figure}
